\documentclass[11pt]{letter}
\usepackage[margin=1in]{geometry}
\usepackage{parskip}

\signature{Hikaru Wakaura \\ Taiki Tanimae \\ QIRI (Quantum Integrated
Research Institute Inc.)}
\address{QIRI (Quantum Integrated Research Institute Inc.)\\
1--16--3 Akasaka, Minato-ku\\ Tokyo 107--0061, Japan\\
h.wakaura@deeptell.jp,\ t.tanimae@deeptell.jp}

\begin{document}
\begin{letter}{Editorial Office \\ New Journal of Physics \\ IOP
Publishing}

\opening{Dear Editors,}

We submit the accompanying manuscript, ``Which quantum algorithms suit
photonic continuous-variable hardware? A comparative resource and
noise-robustness study with a worked spectroscopic demonstration on
H$_2$, HF, and I$_2$,'' for consideration as a Paper in \emph{New
Journal of Physics}. The work falls within the journal's broad scope of
quantum information science, photonic quantum technologies, and
computational molecular physics.

\textbf{Motivation.} Photonic continuous-variable (CV) architectures
support deterministic, room-temperature generation of large cluster
states and are a leading route to fault-tolerant quantum computing,
where overhead is dominated by squeezing and by the consumption of
non-Gaussian resource states. Different algorithms place very different
demands on these resources, yet no unified comparison has quantified
which algorithm classes are intrinsically robust on CV hardware and how
much Gottesman--Kitaev--Preskill (GKP) error correction actually helps.

\textbf{Main results.} We port four representative algorithms---%
control-free phase estimation, randomized product formulas, quantum
Kolmogorov--Arnold networks, and Regev factoring---to a common CV model
and compare them under photon loss (exact Fock-space Kraus channel) and
Gaussian displacement with GKP correction. We establish: (i)~control-%
free phase estimation is the most noise-robust because it needs no
controlled unitaries and, for quadratic Hamiltonians, no non-Gaussian
gates, a property that survives for non-Gaussian Hamiltonians;
(ii)~GKP correction reduces residual error only above the finite-%
squeezing floor (proved), with the gain set by a noise- versus
systematic-limited dichotomy whose empirical envelope constants we
report with $R^{2}\!\ge\!0.89$; (iii)~the comb readout of Regev requires
squeezing scaling as $\sqrt{n}$ (proved); (iv)~against a qubit
surface-code baseline, the CV advantage is structural---Gaussian
operations are free---rather than threshold-based.

\textbf{Worked physics example.} The manuscript includes a quantitative
demonstration on three diatomic molecules (H$_2$, HF, I$_2$) spanning a
decade in vibrational quantum ($\omega_e=214$--$4401$~cm$^{-1}$) and an
order of magnitude in anharmonicity. Using experimentally fixed
Huber--Herzberg Morse parameters (no free parameters), CV control-free
QPE recovers the populated vibrational eigenenergies of all three
molecules to within $0.55$~cm$^{-1}$---below the standard
$1$~cm$^{-1}$ spectroscopic threshold; weakly anharmonic I$_2$ reaches
$\sim\!10^{-2}$~cm$^{-1}$. Reaching the same accuracy with a
discrete-variable QPE requires $\sim 1.5\times10^{10}$ T-gates and
$\sim$~16{,}000 controlled time evolutions, while the CV control-free
protocol needs zero controlled unitaries and zero non-Gaussian gates
for the harmonic part. A Chebyshev-KAN classifier trained on the
recovered eigenstates achieves $100\%$ accuracy in assigning
vibrational quantum numbers under realistic measurement noise.

\textbf{Methodological rigour.} The core GKP-correction surrogate is
validated against an explicit Glancy--Knill error-correction Monte
Carlo to within $0.1\%$, lifting a common methodological concern in
CV-resource analyses. A two-dimensional noise phase diagram and
Hamiltonian ensembles confirm the predictions are not artefacts of
single operating points or single instances. We are explicit about
which results are exact, which are validated model surrogates, and
which are extrapolations.

\textbf{Fit to NJP.} New Journal of Physics has a long tradition of
publishing methodological and computational advances at the
intersection of quantum information and physical applications. We
believe the manuscript fits this profile: it provides a unified
methodology for assessing CV algorithm suitability, together with a
quantitative physics application that yields concrete predictions
experimentalists can test on near-term GKP-encoded photonic processors.
All scripts and numerical results that reproduce every figure and
table accompany the submission.

The manuscript has not been published and is not under consideration
at any other journal; the authors declare no competing interests.

\closing{Sincerely,}

\end{letter}
\end{document}
